\documentclass[11pt]{amsart}
\usepackage{amsmath,amssymb,mathtools,booktabs,hyperref}
\usepackage[margin=1in]{geometry}
\newtheorem{theorem}{Theorem}[section]
\newtheorem{lemma}[theorem]{Lemma}
\newtheorem{proposition}[theorem]{Proposition}
\newtheorem{corollary}[theorem]{Corollary}
\theoremstyle{remark}
\newtheorem{remark}[theorem]{Remark}
\newcommand{\R}{\mathbb R}
\newcommand{\Z}{\mathbb Z}
\newcommand{\norm}[1]{\lVert #1\rVert}
\DeclareMathOperator{\Real}{Re}
\DeclareMathOperator{\ess}{ess}
\newcommand{\T}{T}
\title{Symbol calculus for the certification of CGMY spline trajectories\\
\large Technical note for the revision of ``Residual Certification \dots'' (v8)}
\author{Prepared for D.~Damircheli}
\date{October 2026}

\begin{document}
\maketitle

\noindent
This note replaces the three weakest links of v8 (principal-period
assembly, the unvalidated spectral estimate, and the full-order Riesz solve
without preconditioning) by exact statements that cost nothing to evaluate.
Everything below uses the conventions of v8, Section~2: cardinal splines of
order $p\ge2$, basis $\phi_{h,k}(x)=h^{-1/2}N_p((x-x_{\min})/h-k)$, the
backward symbol
\begin{equation}
 A_\mu(\xi)=-(r-\omega_S)\,i\xi-C\Gamma(-Y)\bigl[(M-i\xi)^Y+(G+i\xi)^Y\bigr]+\lambda_0,
 \qquad \mu=(C,G,M,Y,r),
 \label{eq:symbol}
\end{equation}
and the free index set $k=0,\dots,N-p$ ($N_h=N-p+1$).  All numerical values
quoted were produced with \texttt{symbol\_prototype.py} on the v8 parameter box
and reproduce v8's Table~9 and the positivity audit (minimum
$\lambda_{\min}(S_I)=0.01038$ at the same box vertex).

\section{Exact assembly by periodisation}

Write $\vartheta_k:=\vartheta+2\pi k$ and
$\widehat N_p(\vartheta)=\bigl(\tfrac{\sin(\vartheta/2)}{\vartheta/2}\bigr)^p$
(up to a unimodular factor), so that $|\widehat N_p(\vartheta)|^2\le(2/|\vartheta|)^{2p}$.

\begin{lemma}[Periodised symbol]\label{lem:periodisation}
Let $p\ge2$ and $Y\in(0,2)\setminus\{1\}$.  The exact Galerkin matrix of
$\mathcal A_\mu$ in the basis $\{\phi_{h,k}\}$ is the Toeplitz matrix
$(A_h)_{kk'}=t_{k-k'}$ with
\begin{equation}
 t_m=\frac1{2\pi}\int_{-\pi}^{\pi}a_\mu(\vartheta)e^{i\vartheta m}\,d\vartheta,
 \qquad
 a_\mu(\vartheta):=\sum_{k\in\Z}\bigl|\widehat N_p(\vartheta_k)\bigr|^2A_\mu(\vartheta_k/h),
 \label{eq:periodised}
\end{equation}
where the series converges absolutely and uniformly on $[-\pi,\pi]$, and
$a_\mu$ is real-analytic and $2\pi$-periodic.  Likewise the mass matrix,
the first-derivative matrix and the $H^1$ stiffness matrix $h^2R^{(1)}_h$
have symbols
$m(\vartheta)=\sum_k|\widehat N_p(\vartheta_k)|^2$,
$\tfrac{i}{h}d_1(\vartheta)$ with $d_1(\vartheta)=\sum_k|\widehat N_p(\vartheta_k)|^2\vartheta_k$,
and $g_1(\vartheta)=\sum_k|\widehat N_p(\vartheta_k)|^2\vartheta_k^2$.
For $p=3$, $m(\vartheta)=\tfrac{11}{20}+\tfrac{13}{30}\cos\vartheta+\tfrac1{60}\cos2\vartheta$
and $m_{\min}:=\min_\vartheta m=m(\pi)=2/15$.
\end{lemma}

\begin{proof}
With $\widehat\phi_{h,k}(\xi)=h^{1/2}\widehat N_p(h\xi)e^{-i\xi(x_{\min}+kh)}$ and
Plancherel,
$a_\mu(\phi_{h,k'},\phi_{h,k})=\frac1{2\pi}\int_\R A_\mu(\xi)\,h|\widehat N_p(h\xi)|^2e^{i\xi h(k-k')}d\xi$;
substitute $\vartheta=h\xi$ and split $\R=\bigcup_k[-\pi,\pi)+2\pi k$.  Since
$|A_\mu(\xi)|\lesssim 1+|\xi|^{\max(Y,1)}$ and $|\widehat N_p(\vartheta)|^2\lesssim|\vartheta|^{-2p}$,
the terms are bounded by $|\vartheta_k|^{\max(Y,1)-2p}$, summable because
$2p-\max(Y,1)>1$.  Analyticity follows from that of each term and uniform
convergence of derivatives.  The mass symbol is the Euler--Frobenius
polynomial of $N_{2p}$ (entries $N_6(3+m)=\frac1{120}(1,26,66,26,1)$ for $p=3$).
\end{proof}

\begin{lemma}[Truncation tail]\label{lem:tail}
Let $a_{\mu,K}$ be the series \eqref{eq:periodised} truncated to $|k|\le K$,
and for $q\in[0,2p-1)$ put
$\tau_K(q):=\dfrac{2^{2p+1}\pi^{q-2p}(2K-1)^{\,q-2p+1}}{2(2p-q-1)}$.
Then for all $|\vartheta|\le\pi$,
\begin{equation}
 |a_\mu(\vartheta)-a_{\mu,K}(\vartheta)|
 \le \varepsilon_K(\mu):=\frac{|r-\omega_S|}{h}\,\tau_K(1)
 +2C\Gamma(-Y)h^{-Y}\Bigl(1+\frac{\max(G,M)h}{\pi(2K+1)}\Bigr)^{Y}\tau_K(Y)
 +\lambda_0\,\tau_K(0).
 \label{eq:tail}
\end{equation}
\end{lemma}
\begin{proof}
For $|k|\ge K+1$ and $|\vartheta|\le\pi$, $|\vartheta_k|\ge\pi(2|k|-1)\ge\pi(2K+1)$.
Hence $|\widehat N_p(\vartheta_k)|^2|\vartheta_k|^q\le2^{2p}(\pi(2|k|-1))^{q-2p}$ and
$|\delta+i\vartheta_k|^Y\le|\vartheta_k|^Y(1+\delta/(\pi(2K+1)))^Y$.  Sum over
$|k|\ge K+1$ using $\sum_{k\ge K+1}(2k-1)^{-s}\le\int_K^\infty(2t-1)^{-s}dt$.
\end{proof}

\noindent
Numerically (v8 box, $N=64$, $\mu=(0.5,3,5,1.5,0.05)$): the operator-norm
error of $A_{h,K}$ is $3.2\times10^{-4}$ ($K=5$), $3.2\times10^{-6}$ ($K=20$),
$1.2\times10^{-8}$ ($K=100$), $3.7\times10^{-12}$ ($K=1000$), and the bound
\eqref{eq:tail} is within a factor $3$ of the observed error at $K=5$ and within $1.5$--$1.7$ for $K\ge20$.
The principal-period rule of v8 is the case $K=0$: its relative operator-norm
error is $0.387,0.378,0.369,0.362$ at $N=64,128,256,512$, matching v8's
Table~9, because at $\vartheta=\pm\pi$ the $k=0$ and $k=\mp1$ terms have equal
modulus.  With $K\ge100$ the assembled matrices \emph{are} the exact Galerkin
matrices to $10^{-8}$, and every exact-assembly statement of v8 applies to
them with the explicit perturbation \eqref{eq:tail}.

\section{Toeplitz sections and the free system}

Let $\T_n(f)$ denote the $n\times n$ Toeplitz matrix with symbol $f$.  Because
the free indices are contiguous, the free matrices of v8 are sections:
$A_I=\T_{N_h}(a_\mu)$, $M_I=\T_{N_h}(m)$, $S_I=\T_{N_h}(s_\mu)$ with
$s_\mu:=\Real a_\mu$ (the drift symbol is odd, so it drops out of the
symmetric part), and $h^2R^{(1)}_I=\T_{N_h}(g_1)$, which is banded
(bandwidth $p-1$).

\begin{lemma}[Symbol inequalities]\label{lem:toeplitz}
For bounded $2\pi$-periodic $f,g$ with $g>0$ real and every $v\in\R^n$,
\[
 v^\top\T_n(f)v=\frac1{2\pi}\int_{-\pi}^{\pi}f(\vartheta)\Bigl|\sum_kv_ke^{ik\vartheta}\Bigr|^2d\vartheta,
 \qquad
 \norm{\T_n(f)}_2\le\norm f_\infty,
 \qquad
 v^\top\T_n(f)v\ \ge\ \Bigl(\inf_\vartheta\frac{f}{g}\Bigr)\,v^\top\T_n(g)v .
\]
In particular $\lambda(\T_n(f))\subset[\inf f,\sup f]$ for real $f$.
\end{lemma}
\begin{proof}
Expand the quadratic form; the norm bound is the compression of
multiplication by $f$ on $L^2(\mathbb T)$; the last inequality is the first
identity applied to $f-(\inf f/g)\,g\ge0$.
\end{proof}

\section{Spectral lower bound and the extension to $Y<1$}

\begin{lemma}[L\'evy--Khintchine form of the real symbol]\label{lem:LK}
For $Y\in(0,2)\setminus\{1\}$, $a>0$ and $\xi\in\R$,
\[
 C\Gamma(-Y)\bigl[a^Y-\Real(a+i\xi)^Y\bigr]
 =C\,I_a(\xi),\qquad
 I_a(\xi):=\int_0^\infty(1-\cos\xi x)\,e^{-ax}x^{-1-Y}dx\ \ge0 .
\]
Consequently $\Real A_\mu(\xi)=r+C\,[I_G(\xi)+I_M(\xi)]\ge r$ for all $\xi$,
without any restriction $Y>1$.
\end{lemma}
\begin{proof}
For $Y\in(1,2)$: $\int_0^\infty(e^{-zx}-1+zx)x^{-1-Y}dx=\Gamma(-Y)z^Y$,
$\Real z\ge0$; apply with $z=a+i\xi$ and $z=a$, subtract, use
$\Gamma(1-Y)=-Y\Gamma(-Y)$ and take real parts.  For $Y\in(0,1)$ use
$\int_0^\infty(e^{-zx}-1)x^{-1-Y}dx=\Gamma(-Y)z^Y$ in the same way; the
sign of $\Gamma(-Y)$ cancels against that of the bracket.
\end{proof}

\begin{theorem}[Rigorous spectral lower bound]\label{thm:spectral}
For exact assembly and every $N$,
\begin{equation}
 \lambda_{\min}(S_I)\ \ge\ s_{\min}(\mu):=\min_{\vartheta}s_\mu(\vartheta)
 \ \ge\ r\,m_{\min},
 \qquad S_I\succeq r\,M_I .
 \label{eq:spectral}
\end{equation}
For the truncated assembly with $K$ aliases, subtract $\varepsilon_K(\mu)$ of
\eqref{eq:tail}.  Both bounds are parameter-explicit and valid for
$Y\in(0,2)\setminus\{1\}$.
\end{theorem}
\begin{proof}
$s_\mu=\sum_k|\widehat N_p(\vartheta_k)|^2\Real A_\mu(\vartheta_k/h)\ge r\,m(\vartheta)$
by Lemma~\ref{lem:LK}; apply Lemma~\ref{lem:toeplitz}.
\end{proof}

\noindent
Over 32 box vertices and 256 Sobol points at $N=64$ and $N=256$,
$\lambda_{\min}(S_I)/s_{\min}(\mu)\in[1.003,17.7]$, with the ratio $1.003$
attained exactly at the worst query of v8's audit
($\mu=(0.05,5,9,1.1,0.01)$, $\lambda_{\min}=0.01038$, $s_{\min}=r=0.0100$).
The crude bound $r\,m_{\min}$ is $7.5$--$133\times$ conservative.  $s_{\min}(\mu)$ is
the minimum of an explicit smooth even function of one variable; a validated
enclosure costs one interval evaluation on a grid.  This replaces v8's
Section~4.2 (numerical $\widehat\lambda_{\min}$ estimate) and the
459-coordinate positivity audit by a theorem.  Since the certificate needs
only $S_I\succ0$, \emph{all algebraic results of v8 extend to $Y\in(0,1)$};
only the continuum coercivity in $H^{Y/2}_{M,G}$ keeps the hypothesis $Y>1$.

\begin{remark}[Explicit coercivity constant]
Lemma~\ref{lem:LK} also sharpens v8's Theorem~2.1.  Using $1-\cos u\ge2u^2/\pi^2$
on $|u|\le\pi$,
\[
 I_a(\xi)\ \ge\ \frac{2}{\pi^2}\gamma_Y\,a^{Y-2}\xi^2\ (|\xi|\le a),\qquad
 I_a(\xi)\ \ge\ \frac{2e^{-\pi}}{(2-Y)\pi^Y}\,|\xi|^Y\ (|\xi|\ge a),\qquad
 \gamma_Y:=\int_0^\pi u^{1-Y}e^{-u}du,
\]
which yields $c_1$ of order $10^{-5}$ on the v8 box instead of
$1.6\times10^{-8}$ (Appendix~A of v8), with a two-line proof.
\end{remark}

\section{A banded comparison norm: cheap certificates without Riesz solves}

\begin{theorem}[Comparison norm]\label{thm:comparison}
Fix $\mu$ and define
\[
 \beta_\mu:=\inf_{\vartheta\ne0}\frac{s_\mu(\vartheta)-r\,m(\vartheta)}{g_1(\vartheta)}>0,
 \qquad
 x_\mu:=r\,m+\beta_\mu g_1,\qquad
 X_I:=\T_{N_h}(x_\mu)=r\,M_I+\beta_\mu h^2R^{(1)}_I,
 \qquad
 \Theta_\mu:=\sup_\vartheta\frac{s_\mu(\vartheta)}{x_\mu(\vartheta)}.
\]
$X_I$ is symmetric positive definite and banded with bandwidth $p-1$, and
for every $\rho\in\R^{N_h}$
\begin{equation}
 \rho^\top S_I^{-1}\rho\ \le\ \rho^\top X_I^{-1}\rho\ \le\ \Theta_\mu\,\rho^\top S_I^{-1}\rho .
 \label{eq:comparison}
\end{equation}
\end{theorem}
\begin{proof}
$s_\mu-r\,m=C\sum_k|\widehat N_p(\vartheta_k)|^2[I_G+I_M](\vartheta_k/h)\ge0$ vanishes
only at $\vartheta=0$, where both it and $g_1$ are $\asymp\vartheta^2$
($I_a(\xi)\sim\tfrac12 I_a''(0)\xi^2$ with $I_a''(0)=\Gamma(2-Y)a^{Y-2}>0$), so
$\beta_\mu>0$.  Lemma~\ref{lem:toeplitz} gives $S_I\succeq X_I\succeq\Theta_\mu^{-1}S_I$;
invert in the Loewner order.
\end{proof}

\begin{corollary}[Cheap certificate]\label{cor:cheap}
In v8's Theorem~4.1 and Proposition~4.5 replace
$\norm{\rho^{n+1/2}}_{S_h^{-1}}^2$ by $\rho^\top X_I^{-1}\rho$.  The resulting
$\eta_{\rm aug}^X$ is a rigorous upper bound with effectivity at most
$\sqrt{\Theta_\mu}$ times that of the exact estimator, and costs one banded
Cholesky factorisation plus $N_T$ banded solves ($O(N_TN_h)$ operations).
No Riesz solve, no spectral estimate and no bracket are needed.
\end{corollary}

\noindent
Measured on the v8 box ($\sqrt{\Theta_\mu}$ is the a-priori bound; the last
columns are observed ratios $\norm\rho_{X^{-1}}/\norm\rho_{S^{-1}}$ for
residuals $\rho=A_Ie$ with smooth $e$):
\begin{center}\small
\begin{tabular}{rcccc}
\toprule
$N$ & $\sqrt{\Theta_\mu}$ median & $\sqrt{\Theta_\mu}$ max & observed median & observed max\\
\midrule
64 & 1.15 & 1.74 & 1.09 & 1.56\\
256 & 1.49 & 2.88 & -- & --\\
512 & 1.71 & 3.76 & 1.37 & 2.47\\
1024 & 2.07 & 4.91 & -- & --\\
4096 & 2.89 & 8.27 & -- & --\\
\bottomrule
\end{tabular}
\end{center}
The growth with $N$ is $\sim(\min(G,M)h)^{(Y-2)/2}$ (the tempering scale moves
below the resolution of a degree-$(p-1)$ trigonometric minorant).  Two
remedies, both optional: a wider-band minorant, or the two-stage rule below.

\begin{proposition}[Two-stage rule]\label{prop:twostage}
Let $\eta^X$ be the cheap bound and $\eta^S$ the exact one.  Then
$\eta^S\le\eta^X\le\sqrt{\Theta_\mu}\,\eta^S$ (with the same $\beta_0$).  Hence:
if $\eta^X\le\varepsilon$ accept; if $\eta^X>\sqrt{\Theta_\mu}\,\varepsilon$
reject without further work (the exact bound cannot pass); otherwise compute
$\eta^S$ by the bracket of v8 with $\alpha_{\rm lb}=s_{\min}(\mu)$.
\end{proposition}

\noindent
For the third case the Riesz systems $S_Iv=\rho$ should be solved by CG with
the Strang circulant preconditioner built from the same symbol samples.
Since $s_\mu$ is continuous and positive, the preconditioned spectrum
clusters at $1$ (Chan--Strang), and we observe $7$--$12$ iterations to
$10^{-12}$ for $N=64,\dots,2048$, against $36$--$586$ for plain CG (v8 hit its
$200$-iteration cap at $N=512$).

\section{Affine decomposition of the fractional blocks and rigorous transfer}

Write
\begin{equation}
 a_\mu=-(r-\omega_S)\tfrac{i}{h}d_1+\lambda_0\,m-C\Gamma(-Y)h^{-Y}\Phi_{\nu},
 \qquad
 \Phi_\nu(\vartheta):=\sum_k|\widehat N_p(\vartheta_k)|^2\bigl[(Gh+i\vartheta_k)^Y+(Mh-i\vartheta_k)^Y\bigr],
 \label{eq:affine-split}
\end{equation}
with $\nu=(G,M,Y)$.  The first two terms are already affine in $\mu$ with
fixed matrices $D_h$, $M_h$.  What has to be interpolated is the
\emph{three-parameter family of periodic functions} $\{\Phi_\nu\}$, not the
matrix family: because $\T_n(\cdot)$ is linear, an interpolant
$\Phi_\nu\approx\sum_{q=1}^Q\ell_q(\nu)\Phi_{\nu_q}$ gives
$A_Q(\mu):=-(r-\omega_S)D_h+\lambda_0M_h-C\Gamma(-Y)h^{-Y}\sum_q\ell_q(\nu)\T(\Phi_{\nu_q})$
with $Q+2$ terms and the \emph{same} coefficient functions.

\begin{lemma}[Error transfer]\label{lem:eim}
Let $e_\nu:=\Phi_\nu-\sum_q\ell_q(\nu)\Phi_{\nu_q}$.  Then
\[
 \norm{A_h(\mu)-A_Q(\mu)}_2\ \le\ \varepsilon_Q(\mu):=C\Gamma(-Y)h^{-Y}\norm{e_\nu}_\infty,
\]
and $\norm{e_\nu}_\infty\le\max_{j}|e_\nu(\vartheta_j)|+\tfrac{\pi}{L}\norm{e_\nu'}_\infty$
on an $L$-point grid, with $e_\nu'$ obtained from the explicit derivative
$\partial_\vartheta(\delta+i\vartheta)^Y=iY(\delta+i\vartheta)^{Y-1}$ and the
same tail bound.  Evaluating $\varepsilon_Q(\mu)$ at a new query costs
$O(LK)$ scalar operations and no linear algebra.
\end{lemma}
\begin{proof}
Lemma~\ref{lem:toeplitz}, second inequality, and the mean value theorem.
\end{proof}

\begin{theorem}[Certificate with an affine surrogate operator]\label{thm:eim-cert}
Let the candidate trajectory be produced by any method (e.g.\ the reduced
system $V^\top A_Q(\mu)V$), and let $\rho_Q^{n+1/2}$ be its residual
evaluated with $A_Q$.  Then its residual with respect to the true assembled
system satisfies
\[
 \norm{\rho^{n+1/2}}_{S_I^{-1}}\ \le\ \norm{\rho_Q^{n+1/2}}_{S_I^{-1}}
 +\varepsilon_Q(\mu)\,\alpha_{\rm lb}^{-1/2}\,\norm{\overline c_\theta^{\,n+1/2}}_2,
 \qquad \alpha_{\rm lb}=s_{\min}(\mu),
\]
and v8's Theorem~4.1 holds with $\eta^2=\beta_0^2+\Delta\tau\sum_n(\norm{\rho_Q}_{S^{-1}}+\varepsilon_Q\alpha_{\rm lb}^{-1/2}\norm{\overline c_\theta}_2)^2$.
The same holds with $X_I$ in place of $S_I$ and $\min_\vartheta x_\mu(\vartheta)\ (\ge r\,m_{\min})$ in place of $\alpha_{\rm lb}$.
\end{theorem}
\begin{proof}
$\rho=\rho_Q+(A_h-A_Q)\overline c_\theta$; bound the second term by
$\norm{(A_h-A_Q)\overline c}_2\lambda_{\min}(S_I)^{-1/2}$ and use
Theorem~\ref{thm:spectral}.  This is exactly the allowance structure of v8,
Remark~4.4, with $d_n$ now explicit.
\end{proof}

\noindent
Greedy EIM on $\{\Phi_\nu\}$ over the v8 box (training $12^3$ tensor grid,
test 400 Sobol points, sup-norm error relative to $\max|\Phi|$):
\begin{center}\small
\begin{tabular}{lcccc}
\toprule
 & $10^{-4}$ & $10^{-6}$ & $10^{-8}$ & $10^{-10}$\\
\midrule
$Q$ needed, $N=64$ ($h=1/4$) & 8 & 14 & 18 & 23\\
$Q$ needed, $N=512$ ($h=1/32$) & 10 & 16 & 23 & 30\\
\bottomrule
\end{tabular}
\end{center}
Per block (two parameters $(\delta,Y)$) the numbers are $7$--$8$, $8$--$11$,
$11$--$15$, $13$--$18$.  The singular values of the snapshot matrix decay
geometrically (ratio $<10^{-10}$ at index $21$--$25$), as expected for a family
analytic in $(\delta,Y)$ with the nearest singularity at $\delta=0$, i.e.\
at Bernstein-ellipse distance set by $G_{\max}/G_{\min}$ and $M_{\max}/M_{\min}$
(not by $h$).  Thus $A_h(\mu)$ over the full five-dimensional box is affine
with about $25$ terms to $10^{-8}$.

\begin{remark}[Reduced online certificate]
With a fixed comparison norm $\bar X=\T(\bar x)$ (one per cell of a coarse
partition of the $(r,C)$ range) and the affine $A_Q$, the standard
offline/online decomposition of the residual norm applies: the online cost of
$\eta^{\bar X}$ for a rank-$r$ candidate is $Q^2(2r+1)^2$ operations,
independent of $N_h$ ($\approx6\times10^6$ for $Q=25,r=48$, $50$~MB offline
storage), and
$\rho^\top S_I^{-1}\rho\le\bigl(\inf_\vartheta s_\mu/\bar x\bigr)^{-1}\rho^\top\bar X^{-1}\rho$
by Lemma~\ref{lem:toeplitz}, the infimum being again an explicit
one-dimensional computation at the query.
\end{remark}

\section{Where the framework pays: a certificate for the American put}

For a discrete obstacle problem in coefficient space (B-splines are
nonnegative, so $c\ge\psi_h$ coefficientwise implies $U_h\ge\psi_h$), one
backward-Euler step reads: find $u^{n+1}\ge\psi_h$ with
$(B_+u^{n+1}-b^n)^\top(v-u^{n+1})\ge0$ for all $v\ge\psi_h$, where
$B_+=M_h+\Delta\tau A_h$ and $b^n=M_hu^n+\Delta\tau f^{n+1}$.

\begin{theorem}[Residual certificate for the discrete variational inequality]
\label{thm:vi}
Let $\{u_\theta^n\}$ be any candidate with $u_\theta^n\ge\psi_h$ for all $n$,
let $R^n:=B_+u_\theta^{n+1}-M_hu_\theta^n-\Delta\tau f^{n+1}$, and split
$R^n=R^n_+-R^n_-$ into positive and negative parts.  Then for every $m$,
\[
 \norm{e^m}_{M_h}^2+\Delta\tau\sum_{n<m}\norm{e^{n+1}}_{S_h}^2
 \ \le\ \norm{e^0}_{M_h}^2
 +\sum_{n<m}\Bigl[\,2\,(R^n_+)^\top(u_\theta^{n+1}-\psi_h)
 +\Delta\tau^{-1}\norm{R^n_-}_{S_h^{-1}}^2\Bigr].
\]
\end{theorem}
\begin{proof}
Test the inequality with $v=u_\theta^{n+1}$:
$(B_+e^{n+1}-M_he^n,\,e^{n+1})\le-(R^n,e^{n+1})=(R^n,u_\theta^{n+1}-u^{n+1})$.
Since $u^{n+1}\ge\psi_h$ and $R_+\ge0$, $(R^n_+,u_\theta^{n+1}-u^{n+1})\le(R^n_+,u_\theta^{n+1}-\psi_h)$;
and $(R^n_-,e^{n+1})\le\norm{R^n_-}_{S^{-1}}\norm{e^{n+1}}_S$.  The left side
equals $\tfrac12(\norm{e^{n+1}}_M^2-\norm{e^n}_M^2+\norm{e^{n+1}-e^n}_M^2)+\Delta\tau\norm{e^{n+1}}_S^2$;
Young's inequality and summation finish the proof.
\end{proof}

\noindent
Everything in Sections~1--5 (exact assembly, $s_{\min}$, the banded norm,
the affine $A_Q$) applies verbatim; only the right-hand side gains the
complementarity term, which vanishes where the candidate is either off the
obstacle with nonpositive residual or on it.  The Crank--Nicolson version
requires testing the inequality at the midpoint, which is not admissible in
general; this is the one open technical point and the reason to use backward
Euler (or Rannacher-type restarts) for the American reference.

\end{document}
